\section{Conclusion}
\label{sec:conclusion}

We gave an EEG tokenizer the one relation its tokens never carry, the
coupling between frequency bands, and measured what a clinical model
gains from having it explicitly. In a small encoder built around the
token, removing the coupling content costs between $0.03$ and $0.14$ on
every corpus on which we measure it, and the model, at 2.7M parameters,
leads the strongest reproduced foundation-model baselines by $0.15$ and
$0.09$ AUC-PR on seizure detection while matching them on event
classification and recording-level abnormality. In place of the native
tokenizer, the token lifts two published encoders trained from scratch
by $0.15$ and $0.24$ $\kappa$ on epileptiform event classification; added
next to an encoder's own spectral branch, the gains are smaller and
corpus-dependent. Where the model loses --- sleep staging, cohort
diagnosis, artifacts, and the smallest seizure corpus --- it loses to
pretrained or larger models. The benefit of explicit cross-frequency
content therefore depends on the encoder's inductive bias, and a
tokenizer claim is only as broad as the set of encoders on which it is
tested; we have reported ours across four. What this work does not yet
establish is the physiological reading of the statistic the token
carries: the coupling is measured at patch resolution over learned bands,
and separating genuine cross-frequency interaction from co-varying
non-stationarity, as well as reporting event-level sensitivity and
false-alarm rates for the seizure corpora, is left to future work.
